%%%
\subsection{Automatyzacja zarządzania zasobami w chmurze}

Automatyzacja zarządzania zasobami w chmurze obliczeniowej jest fundamentalnym elementem wydajnych, 
niezawodnych i skalowalnych systemów informatycznych. Opiera się na wdrażaniu zasad automatyki i 
inżynierii sterowania, umożliwiając inteligentne, często autonomiczne zarządzanie infrastrukturą IT. 
W tym kontekście kluczowym pojęciem pozostaje pętla sterowania (control loop), będąca podstawą systemów 
samonaprawiających się oraz adaptacyjnych rozwiązań chmurowych.

\subsubsection{Pojęcie pętli sterowania (Control Loops) w systemach rozproszonych}

\textbf{Pętla sterowania} to podstawowy mechanizm automatyki, polegający na nieustannym monitorowaniu 
parametrów systemu, ich analizie oraz podejmowaniu działań korygujących, jeśli wykryta zostanie 
odchyłka od zadanych wartości. Pętle te mogą być zarówno otwarte, jak i zamknięte („feedback loops”), 
jednak w praktyce zarządzania chmurą dominują warianty z informacją zwrotną .

Historycznie, idea pętli sterowania wywodzi się z XVIII wieku, gdy James Watt opracował regulator odśrodkowy
 dla maszyn parowych, umożliwiając automatyczne utrzymywanie zadanej prędkości obrotowej silnika poprzez 
 mechaniczne sprzężenie zwrotne.


W architekturach rozproszonych, jak chmura obliczeniowa czy systemy DCS (Distributed Control Systems), 
pętle sterowania realizowane są na wielu poziomach – od lokalnych kontrolerów odpowiedzialnych za pojedyncze węzły, po nadzorczy poziom systemowy. Przykłady obejmują:

\begin{itemize}
    \item \textit{Load Balancing Feedback Loops}: dynamiczne równoważenie obciążenia na podstawie bieżących informacji o stanie zasobów;
    \item \textit{Scaling Feedback Loops}: automatyczna zmiana liczby aktywnych instancji usług zależnie od aktualnych potrzeb;
    \item \textit{Health Monitoring Feedback Loops}: wykrywanie oraz samoczynna reakcja na awarie lub degradację wydajności.
\end{itemize}

Pętle te zapewniają adaptacyjność oraz samonaprawialność systemów, pozwalając na minimalizowanie wpływu awarii,
 optymalizację zużycia zasobów oraz ciągłość działania usług.

\subsubsection{Autonomiczne systemy i samoadaptacja}
Ewolucja automatyki i inżynierii sterowania doprowadziła do powstania \textbf{autonomicznych systemów}, które
nie tylko wykonują zaprogramowane wcześniej zadania, lecz potrafią adaptować swoje zachowanie do zmieniających się 
warunków środowiska i własnego stanu. Zasadniczym mechanizmem jest tu samoadaptacja (self-adaptation), polegająca na dynamicznej zmianie parametrów lub sposobu działania systemu na podstawie analizy bieżących danych.

W systemach autonomicznych typowe są zaawansowane architektury warstwowe, gdzie warstwa zarządzająca (managing subsystem) monitoruje „kondycję” systemu oraz środowiska i na tej podstawie podejmuje decyzje dotyczące rekonfiguracji czy optymalizacji:

\begin{itemize}
    \item Przykładem mogą być pojazdy autonomiczne (np. podwodne roboty inspekcyjne), których sterownik może dynamicznie przełączać tryby działania – np. powrót do stacji ładowania przy niskim poziomie energii, czy wybór mniej energochłonnego algorytmu wykrywania przeszkód w przypadku rosnących trudności środowiskowych ;
    \item W chmurze obliczeniowej takie mechanizmy wykorzystywane są do dynamicznej alokacji zasobów, rekonfiguracji po awariach, czy inteligentnego dostrajania parametrów usług w odpowiedzi na nieprzewidywalne zmiany obciążenia .
\end{itemize}

Badania naukowe wskazują, że skuteczność samoadaptacji i autonomiczności w systemach rozproszonych w dużej mierze zależy od efektywności sprzężenia zwrotnego, precyzyjnego modelowania stanu systemu oraz algorytmów decyzyjnych uwzględniających zarówno aspekty wydajnościowe, jak i bezpieczeństwo czy energooszczędność.

Rozwój autonomicznych i samoadaptacyjnych systemów w chmurze znajduje potwierdzenie w licznych badaniach, zarówno w obszarze zastosowań przemysłowych, jak i rozwoju koncepcji architektonicznych oraz formalnych metod modelowania, analiz i walidacji poprawności działania tych systemów. To właśnie te interdyscyplinarne innowacje – czerpiące z tradycyjnej automatyki, inżynierii oprogramowania, robotyki czy informatyki stosowanej – umożliwiają konsekwentną automatyzację i dalszy rozwój chmury obliczeniowej.


